% Generated draft: requires mathematical and source review.
\documentclass[11pt]{article}
\usepackage[margin=25mm]{geometry}
\usepackage{fontspec}
\setmainfont{Latin Modern Roman}
\usepackage{amsmath,amssymb,mathtools,unicode-math}
\setmathfont{Latin Modern Math}
\usepackage{enumitem,needspace,xurl,hyperref}
\setlength{\parindent}{0pt}
\setlength{\parskip}{5pt}
\setlength{\emergencystretch}{3em}
\newcommand{\N}{\mathbb{N}}
\newcommand{\Z}{\mathbb{Z}}
\newcommand{\Q}{\mathbb{Q}}
\newcommand{\R}{\mathbb{R}}
\newcommand{\C}{\mathbb{C}}
\newcommand{\F}{\mathbb{F}}
\newcommand{\A}{\mathbb{A}}
\newcommand{\PP}{\mathbb{P}}
\newcommand{\E}{\mathbb{E}}
\newcommand{\eps}{\varepsilon}
\newcommand{\dd}{\,\mathrm{d}}
\begin{document}
\section*{UnsolvedMath ID: 30003190}
\textbf{Source problem code:} catalogue entry
\subsection{Definitions and mathematical statement}
For a power series \(C(X)=\sum_{n\geq 0}a_nX^n\) over a valued coefficient ring with valuation \(v\), its Newton polygon is the lower convex hull of the points \((n,v(a_n))\) with \(a_n\neq 0\). Its breaking points are the vertices where the slope changes, together with any endpoints included by the chosen convention. Write \(\operatorname{Br}(N)\) for the breaking points of a Newton polygon \(N\). For the displayed specialization formula, use the usual normalization \(\operatorname{val}_p(p)=1\) and \(|x|_p=p^{-\operatorname{val}_p(x)}\).

\textbf{Classical comparison assertion supplied by the record.} Let \(N_{\mathrm{red}}\) denote the Newton polygon described in the source as that of the characteristic series modulo \(p\), and write
\[
\operatorname{Br}(N_{\mathrm{red}})=\{(n_i,m_i):i\in I\}.
\]
Let \(N_\kappa\) denote the Newton polygon after specialization at a weight \(\kappa\). With the source's weight coordinate \(T\), interpreted through \(T(\kappa)=\kappa(1+p)-1\) wherever this expression is defined, the comparison assertion is
\[
\exists r\in\Q,\quad 0<r<1,\quad
\forall\kappa\text{ with }r<|T(\kappa)|_p<1,\qquad
\operatorname{Br}(N_\kappa)
=
\left\{\left(n_i,\operatorname{val}_p\bigl(\kappa(1+p)-1\bigr)m_i\right):i\in I\right\}.
\]
The annulus has strict inequalities at both boundaries, and the same \(r\) must work for every weight in that annulus.

\textbf{Requested formulation.} Determine the appropriate analogue of this breaking-point assertion for Hilbert modular forms. The original source question identifies this as the Hilbert analogue of part (2) of the Spectral Halo conjecture. It asks what the analogue should be; it does not supply a definite Hilbert conjecture to prove.

\textbf{Unresolved qualifications.} The record does not specify the characteristic series, its operator, its coefficient ring, or the valuation used to form its Newton polygon after reduction modulo \(p\). Reduction modulo \(p\) alone does not specify the relevant Newton polygon; in particular, the meaning of the ordinates \(m_i\) requires clarification. The domain and component of classical weight space, restrictions on \(p\), and the breaking-point endpoint convention are also unstated. For the Hilbert setting, no totally real field, level, local conditions at primes above \(p\), weight coordinates, boundary region, or choice of operator or operators is supplied. Consequently, the record does not determine a unique Hilbert replacement for the annulus or the displayed scaling factor. These choices must remain unresolved rather than being inserted as hypotheses or conclusions.
\subsection{Short English statement}
The problem asks how to formulate the Hilbert modular analogue of a classical spectral-halo assertion: throughout one outer annulus of weight space, the Newton polygon has fixed horizontal breaking-point coordinates and vertically rescaled heights. The supplied record leaves the Hilbert parameters, boundary region, operators, and scaling rule unspecified, so this draft preserves the formulation problem without asserting a particular Hilbert conjecture.
\subsection{Sources}
\begin{itemize}
\item Supplied local UnsolvedMath registry record, permanent ID \(30003190\), title \emph{Spectral Halo Breaking Points for Hilbert Modular Forms}, exact source problem code \texttt{catalogue entry}, in the Oberwolfach Reports Open Problems collection. The record gives proposed year 2016 and names no proposer.
\item Supplied report citation: \emph{Arithmetic Geometry (2017)}, Oberwolfach Reports, pp.\ 2171--2224, DOI \texttt{10.4171/owr/2016/38}. \url{https://doi.org/10.4171/owr/2016/38}. The supplied background describes these reports as records of mathematical workshops and their open-problem discussions. The citation's 2017 label and the 2016 DOI/proposed-year information are retained without reconciliation.
\item Supplied literature source: Ren--Zhao, \emph{Hilbert spectral halo}. \url{https://arxiv.org/abs/2005.14267}.
\item Supplied literature source: Andreatta--Iovita--Pilloni. \url{https://doi.org/10.24033/asens.2362}.
\item Supplied literature source: \emph{Oberwolfach Report 38}. \url{https://doi.org/10.4171/owr/2016/38}.
\item Literature assessment dated exactly \texttt{2026-08-21}: the supplied curation labels the problem open and states that quaternionic spectral-halo results and the Hilbert first construction are established, while the precise Hilbert analogue of the second conjecture remains unformulated or unproved. This is supplied dated curation derived from the OWR record, not an independently verified literature assessment. The meaning of the Hilbert first construction is not explained in the record.
\item The registry labels the statement \texttt{corrected\_verified} and reports that the official report was checked and the reference to part (2) was replaced with the full classical Newton-polygon assertion. This is supplied verification history; this draft does not independently verify the report, its reconstruction, or the mathematical claims.
\end{itemize}
\end{document}
